\begin{algorithm}[H]
\SetAlgoLined
\setstretch{0.92}
\SetAlFnt{\normalsize}
\SetArgSty{textnormal}
\SetKwInput{KwIn}{Input}
\SetKwInput{KwInit}{Initialize}
\SetKwInput{KwOut}{Output}
\SetKwFor{For}{for}{do}{end for}
\SetKwFor{ForEach}{for each}{do}{end for}
\SetKwIF{If}{ElseIf}{Else}{if}{then}{else if}{else}{end if}
\SetKw{KwRet}{return}
\SetKw{KwContinue}{continue}
\caption{SFETA：完整合法域动态边际任务分配（迭代收敛版）}
\label{alg:q2-sfeta}
\KwIn{任务集 $\mathcal I$，完整能源基准状态，网络/资源边界，目标模式 $m\in\{\mathrm{Cost},\mathrm{Carbon}\}$，相对收敛容差 $\epsilon$，最大完整扫描次数 $K_{\max}$}
\KwInit{由附件 source-local immediate 状态构造 $\mathbf x^{(0)}$；若该状态存在 GPU 超限，则先在完整合法域内执行可行性修复；$J^{(0)}\leftarrow+\infty$}
\BlankLine
\ForEach{任务 $i\in\mathcal I$}{
    构造空间合法域 $R_i=\{r:\mathrm{Latency}[o_i,r]\le L_i\}$\;
    构造时间合法域 $T_i$：RT 仅含 $a_i$；Batch/Training 使用全部满足 Arrival、LatestFinish 与 2406 的整数开工时刻\;
    记 $\Omega_i=R_i\times T_i$；若 $\Omega_i=\varnothing$ 则返回 FAIL\;
}
按 $|R_i|$ 升序 $\to$ Slack 升序 $\to$ GPU-hour 降序固定任务扫描顺序\;
\BlankLine
\For{$k\leftarrow1$ \KwTo $K_{\max}$}{
    从上一遍排程 $\mathbf x^{(k-1)}$ 重构 GPU/IT/Facility 与逐时能源状态\;
    $\mathrm{changed}\leftarrow0$，$\mathrm{fallback}\leftarrow0$\;
    \ForEach{排序后的任务 $i\in\mathcal I$}{
        读取任务 $i$ 的当前 placement $(r_i^{old},s_i^{old})$\;
        \textbf{暂时移除任务 $i$}，得到不含该任务的残余 workload 状态与能源状态 $S^{-i}$\;
        \ForEach{候选 $(r,s)\in\Omega_i$}{
            将任务 $i$ 按 fractional overlap 加回 $S^{-i}$，形成候选状态 $S^{-i}\oplus(i,r,s)$\;
            检查 GPU-hour、IT、Facility、GridImport、SLA、LatestFinish、finish$\le2406$；RT 额外检查 $s=a_i$\;
            若移除旧任务使某些能源小时暂时触及非负边界，则仅保留能够覆盖并修复这些 removal-dependent 小时的候选\;
            对通过硬约束的候选，计算动态边际响应
            \[
            \Delta E_{i,r,s}=\mathrm{EnergyResponse}(S^{-i}\oplus(i,r,s))-\mathrm{EnergyResponse}(S^{-i}),
            \]
            并由其得到 $\Delta\mathrm{GridPurchase}$、$\Delta\mathrm{Curtailment}$、$\Delta\mathrm{Cost}$ 与 $\Delta\mathrm{Carbon}$\;
        }
        得到全部合法候选集合 $\mathcal F_i$\;
        \If{$\mathcal F_i=\varnothing$}{
            对旧 placement 重新执行完整硬约束检查；若仍不可行则返回 FAIL，否则保留旧 placement 并令 $\mathrm{fallback}\leftarrow\mathrm{fallback}+1$\;
        }
        \Else{
            \If{$m=\mathrm{Cost}$}{按 $\Delta\mathrm{Cost}$ 升序、$\Delta\mathrm{Carbon}$ 次序选择候选\;}
            \Else{按 $\Delta\mathrm{Carbon}$ 升序、$\Delta\mathrm{Cost}$ 次序选择候选\;}
            对数值等价候选执行 tie-break：更早开工 $\to$ 保持 SourceRegion $\to$ 更低网络时延\;
            将最优候选重新插入当前状态；若 placement 改变则 $\mathrm{changed}\leftarrow\mathrm{changed}+1$\;
        }
    }
    由完整排程重新计算真实目标 $J^{(k)}$、硬约束审计和能源平衡残差\;
    \If{$k\ge2$ 且 $(J^{(k-1)}-J^{(k)})/|J^{(k-1)}|<\epsilon$ 且 $\mathrm{fallback}=0$ 且硬约束全部 PASS}{
        \textbf{break}\;
    }
}
\BlankLine
计算 $\mathrm{Cost}$、$\mathrm{Carbon}$、$\eta_R$、迁移/等待/网络时延及收敛轨迹；对 Cost-primary 可从不同初始可行排程重复运行用于路径依赖审计\;
\KwRet 调度方案 $\mathbf x$、主目标、平行评价指标和硬约束审计\;
\KwOut{各任务执行区域与开工时刻、Cost/Carbon/新能源利用率、迁移率、等待分布、收敛与初始化稳健性结果}
\end{algorithm}